\section{Introduction}
\label{sec:intro}

A weighted average of several classifiers' predicted probabilities is among the oldest and most reliable devices in
applied machine learning \citep{hansen1990neural,wolpert1992stacked,kittler1998combining,caruana2004ensemble}. Its
weights are not free parameters: to keep the output a probability they must be non-negative and sum to one, and
non-negativity with normalization is what makes stacked combinations stable \citep{breiman1996stacked}. The feasible
set is the unit simplex, so ensemble weighting is a \emph{compositional} optimization problem.

That has a methodological consequence: optimization over a simplex is the subject of mixture experiments, a mature
branch of experimental design \citep{scheffe1958mixtures,scheffe1963simplexcentroid,cornell2002mixtures}. The
classifiers are the components, a pure vertex is a single model, the centroid uniform soft voting, and a Scheffé
polynomial the natural response surface. The correspondence is exact and already exploited for portfolio weights
\citep{mendes2016portfolio_mde,leal2022portfolio_doptimal_mde} and neural-network ensembles
\citep{moreira2021ann_ensemble_mde}, and combined with Normal Boundary Intersection for forecast combination
\citep{bacci2019nbi_mixture}, load forecasting \citep{rocha2025ensemblestlf} and multiobjective engineering
optimization \citep{pereira2025hybrid}. Most directly, \citet{kwon2024ensemble} treat five heterogeneous
\emph{classifiers}' weights as mixture components and maximize a Scheffé polynomial of accuracy or $F_1$ on eleven
tabular binary datasets.

\textbf{We therefore claim no part of that construction:} mixture design over combination weights, a Scheffé surrogate
over the simplex, and NBI on that surrogate are established prior work, individually and for classifier ensembles.
What is not established is whether the construction stays \emph{trustworthy} once the problem becomes genuinely
multiobjective --- this paper's question.

\subsection{Why the transfer is not routine}

The applications in which this framework matured share three properties that classifier ensembles do not.

First, \emph{the responses are smooth}: tool life or process cost varies smoothly with cutting speed, where a
low-order polynomial is a defensible local model. The area under the ROC curve is instead a rank statistic, piecewise
constant in the weights. No experimental design repairs a misspecified model class, and a smooth surrogate fitted to a
step-like response misrepresents it systematically.

Second, \emph{the decision variables are the mixture}. When the simplex holds scalarization weights, as in the lineage
above, NBI's anchors are optima of the surrogate in another space; when it holds the decision variables themselves,
the anchors are feasible-region points --- often vertices --- so an anchor error displaces the geometry in which every
subproblem is posed. NBI derives its whole search geometry from the payoff matrix of individual optima
\citep{das1998nbi}, so errors there propagate globally. That misplaced anchors damage the construction is established:
payoff-table minima can differ substantially from the true minimum criterion values over the efficient set
\citep{isermann1988payoff}, the nadir point is hard to estimate \citep{deb2010nadir}, and the choice of the individual
minima has itself been treated as a design lever \citep{herrmann2026nonextreme}. What that literature, and the
normalization remedies built on it \citep{messac2003normalized}, asks about is estimating the \emph{true} objectives;
none of it is about a surrogate standing in for them.

Third, \emph{one objective is a step function}: a model with any non-negligible weight must be loaded and executed,
one with none costs nothing, so deployment cost is not proportional to a weight. A method requiring differentiable
objectives can only optimize a continuous relaxation, with no guarantee that it and the quantity actually paid agree
about which solution is best.

Each breaks a different part of the pipeline --- the model class, the search geometry, the objective definition ---
and none is visible in the aggregate fit statistics such studies conventionally report.

\subsection{Approach}

We decompose the framework. Because the ensemble output is linear in the weights and base-model probabilities can be
cached, evaluating a weight vector on the true objectives costs one matrix--vector product ($0.4$ to
$2.1$~milliseconds here) --- cheap enough to run the \emph{same} NBI three times per replication, changing only where
objectives and anchors come from.

\begin{itemize}
\item \textbf{NBI-A} optimizes the Scheffé surfaces with anchors from their own optima: the inherited construction and
our control.
\item \textbf{NBI-B} optimizes the same surfaces with anchors from the real single-objective optima; A versus B
isolates \emph{anchor misplacement}.
\item \textbf{NBI-C} drops the metamodel and runs the same NBI on the exact out-of-fold objectives --- a
metamodel-free NBI in the published sense of \citet{gellerich2023doenbi}, whose subproblems are physical experiments;
B versus C isolates residual \emph{interior surrogate error}.
\end{itemize}

All three arms fit their surfaces once from a fixed design and then optimize. This is the \emph{sequential} form of
surrogate-assisted optimization \citep{tabatabaei2015survey} --- one-shot, with no infill --- and it is the form the
lineage above uses; nothing here tests model-managed, infill-based surrogate-assisted optimizers such as ParEGO,
MOEA/D-EGO or SMS-EGO \citep{knowles2006parego,zhang2010moeadego,ponweiser2008smsego}, which evaluate new points on
the true objectives inside the search loop and exist precisely to repair the kind of surrogate error diagnosed here.
Every conclusion in this paper is scoped to the one-shot, no-infill case (\cref{sec:rel:surrogate}).

We also run NSGA-II on the same objectives at a matched real-objective evaluation budget. Four controls the lineage
does not contain surround the decomposition: external validation of every surface on 100 unseen compositions, with a
reliability gate fixed before the replicated benchmark; revalidation of every candidate on the true objectives before
any indicator; an empirical Pareto reference whose sampled core of at least $10^{5}$ points per replication is built
independently of the surrogate and checked for convergence, then augmented with every candidate set for scoring; and
an untouched holdout never used for fitting, weighting, thresholding or selection.

Four public binary tabular datasets --- customer-transaction prediction (Santander), insurance-claim processing (BNP
Paribas Cardif), auto-insurance claim prediction (Porto Seguro) and consumer-credit default (UCI credit) --- are
partitioned 30 times each, giving 120 replications with recorded seeds. The four prediction tasks differ, but all four
come from financial-services applications, so any cross-dataset statement is made within a single application domain.
Because partitions of one dataset overlap heavily, we treat the replication distribution as partition sensitivity, not
a sample from a population: comparisons are paired by partition and tested with the Nadeau--Bengio correction
\citep{nadeau2003inference}, the dataset is the unit of generalization, and nothing is pooled across the 120
replications.

\subsection{Research questions}

\begin{enumerate}
\item[\textbf{RQ1}] How accurately do Scheffé surfaces predict ensemble performance on compositions they have not
seen, and does this depend on the dataset, on the metric, or on both?
\item[\textbf{RQ2}] How stable are the fitted surfaces across resampling, and does the classical mixture-design
reading of the interaction coefficients survive confrontation with the real objectives?
\item[\textbf{RQ3}] What goes wrong when NBI is anchored on surrogate optima?
\item[\textbf{RQ4}] Does replacing surrogate anchors with real ones repair the front, and does the reliability gate
predict when this is necessary?
\item[\textbf{RQ5}] How does surrogate-assisted NBI compare with metamodel-free NBI, what does the difference cost,
and how do both compare with a canonical evolutionary optimizer at a matched evaluation budget?
\item[\textbf{RQ6}] How large is the actual conflict between discrimination and calibration?
\item[\textbf{RQ7}] Does the definition of deployment cost change which method wins?
\item[\textbf{RQ8}] Do solutions selected on out-of-fold data transfer to an untouched holdout?
\item[\textbf{RQ9}] Would a third of the replications have produced the same conclusions?
\end{enumerate}

\subsection{Contributions}

\begin{enumerate}
\item \textbf{A controlled decomposition of a one-shot surrogate-assisted multiobjective pipeline into surrogate
error, anchor misplacement and solver behaviour}, through three NBI variants defined by two methodological factors.
The ingredients are not new and are not claimed: anchor and ideal/nadir misestimation is an established failure mode
\citep{isermann1988payoff,deb2010nadir,he2021normalization}, the choice of individual minima has already been treated
as a design lever \citep{herrmann2026nonextreme}, and NBI run metamodel-free on measured objectives with real anchors
is published \citep{gellerich2023doenbi}. What we did not find is a study that computes the anchors of a CHIM method
both from a fitted surrogate and from the true objectives on the same problem and compares the resulting fronts, so
what is new here is the surrogate \emph{provenance} of the anchor error and the controlled comparison of the two
constructions; the nearest ensemble-weighting precedent \citep{kwon2024ensemble} optimizes one metric at a time and
has no anchors or front.

\item \textbf{Evidence that, among failures attributable to the surrogate pipeline, anchor misplacement has the
largest observed effect.} Real anchors improve both primary endpoints in 30 of 30 partitions on the two datasets whose
ROC-AUC surfaces almost never pass the gate, in 23 (IGD$^+$) and 24 (hypervolume) of 30 on Porto Seguro, and in 24 of
30 on UCI credit. Removing the surrogate entirely, once anchors are correct, adds much less --- except on UCI credit,
where subproblem convergence rather than surrogate error dominates, the counterexample that prevents a universal
claim.

\item \textbf{A demonstration that an external reliability gate detects unusable surfaces but not misplaced anchors.}
On one dataset the two properties diverge: the partitions where surrogate-anchored NBI collapses are precisely those
where the parsimony rule selects a surface that fits held-out compositions \emph{better} and passes the gate more
often. Surface quality and argmax quality are different properties, and only the second governs NBI.

\item \textbf{A replicated evaluation architecture for ensemble-weight optimization} --- external surrogate
validation, real-objective revalidation, a reference with a method-independent sampled core and holdout confirmation,
over 4 datasets $\times$ 30 partitions with overlap-corrected paired inference. Three findings it overturned --- two
of them artifacts that survived an earlier ten-partition analysis --- are reported explicitly.

\item \textbf{Two findings about how these surfaces should be read.} The classical synergism criterion $\beta_{ij} >
|\beta_i - \beta_j|$ holds for the leading pair's fitted surfaces on all four datasets in every partition, yet the
real blend beats its better member on only one; across all pairs the surrogate over-predicts blending gains 20 times
and under-predicts none. And the continuous cost these methods optimize disagrees with the step cost deployment pays
often enough to change the winning method in up to 24 of 30 partitions. That a change of cost accounting can reorder
a comparison is known across ensemble and cascade architectures \citep{wang2022committees}; what is reported here is
the disagreement between two accountings applied to one and the same weight vector.
\end{enumerate}

We claim neither that NBI outperforms other multiobjective optimizers --- at a matched evaluation budget NSGA-II
outperforms its metamodel-free arm on all four datasets --- nor that it produces more uniformly spaced fronts, nor
that any method here is a general-purpose recommendation. The contribution is the set of conditions under which an
established framework can and cannot be trusted in a new problem class, and the controls needed to tell the
difference --- stated for the one-shot, no-infill form of that framework, since no arm here performs iterative infill.
